% ============================================================================
%  Activity cliffs on human acetylcholinesterase: a pre-registered escalation
%  from 2D descriptors to molecular dynamics
%
%  Build:  latexmk -pdf main.tex        (latexmk and bibtex are installed;
%                                        biber is NOT, so this uses natbib)
%  Figures come from ../figures, generated tables from ../results.
%  Every number in the text should trace to a CSV in ../results -- see
%  ../results/INDEX.md for which stage produced what.
% ============================================================================
\documentclass[11pt,a4paper]{article}

\usepackage[utf8]{inputenc}
\usepackage[T1]{fontenc}
\usepackage[margin=2.5cm]{geometry}
\usepackage{graphicx}
\usepackage{array}
\usepackage{booktabs}
\usepackage{amsmath}
\usepackage{siunitx}
\usepackage[round]{natbib}
\usepackage{xcolor}
\usepackage[colorlinks=true,allcolors=blue]{hyperref}
\usepackage{lineno}

\graphicspath{{../figures/}}

% Drafting aids. Comment out \draftingtrue before submission and both of these
% disappear; \TODO is deliberately loud so nothing ships half-written.
\newif\ifdrafting \draftingtrue
\ifdrafting
  \newcommand{\TODO}[1]{\textcolor{red}{\textbf{[TODO: #1]}}}
  \newcommand{\NOTE}[1]{\textcolor{teal}{\small[#1]}}
  \linenumbers
\else
  \newcommand{\TODO}[1]{}
  \newcommand{\NOTE}[1]{}
\fi

% Shorthands for the three framings, used constantly and easy to get wrong.
\newcommand{\framingA}{framing~A}
\newcommand{\framingB}{framing~B}
\newcommand{\framingC}{framing~C}

\title{Activity cliffs on human acetylcholinesterase:\\
       a pre-registered escalation from 2D descriptors to molecular dynamics}
\author{%
  Isman Kurniawan\thanks{Corresponding author.
    \texttt{isman.krn@gmail.com}. ORCID \href{https://orcid.org/0000-0003-3485-3063}%
    {0000-0003-3485-3063}.} \\[2pt]
  \normalsize School of Computing, Telkom University%
}
\date{\today}

\begin{document}
\maketitle

\input{sections/00\_abstract}
\input{sections/01\_introduction}
\input{sections/02\_methods}
\input{sections/03\_results}
\input{sections/05\_conclusions}

\section*{Data and code availability}
All analysis code is in a public repository at
\url{https://github.com/ismankrn/qsar-md-ache-cliffs}, archived at
\url{https://doi.org/10.5281/zenodo.23498751}. The deposit carries the 229 per-run descriptor files, one JSON per
trajectory and \SI{1.8}{\mega\byte} in total, together with the curated ligand set, the
frozen splits and every derived table the figures and the reported numbers are computed
from. Those files are sufficient to reproduce every result in this paper without
rerunning a simulation.

The trajectories themselves are not deposited. Reduced protein and ligand coordinates at
\SI{100}{\pico\second}, \SI{2.6}{\giga\byte} over the 229 runs, are available from the
corresponding author on request. Any trajectory supplied this way needs the treatment of
the periodic boundary described in Section~\ref{sec:md} before anything is measured on
it, since the archives are wrapped molecule-wise and the capped receptor is three
separate chains. Over the 229 archives an RMSD computed directly on one returns a
median of \SI{45}{\angstrom} for the receptor backbone and \SI{24}{\angstrom} for the
ligand, against the \SI{1.4}{\angstrom} and \SI{3.0}{\angstrom} the same quantities
give once the boundary has been handled.

\section*{Acknowledgements}
The author thanks QuasiLab and Mahameru BRIN \citep{mahameru2025} for their
mini-cluster and HPC facilities.

\bibliographystyle{plainnat}
\bibliography{refs}

\input{sections/90\_si}

\end{document}
